% !TEX root = ../main.tex
\chapter{BIC / quasi-BIC、\texorpdfstring{$\Gamma$}{Gamma} 点辐射与远场工程}
\label{ch:bic}

\paragraph{读完本章，你应能}
\begin{itemize}
  \item 理解连续谱中的束缚态（bound state in the continuum, BIC）与准 BIC 在 PCSEL 中为什么重要。
  \item 学会把 $\Gamma$ 点、辐射损耗、偏振、远场工程和 BIC 语言明确接起来，而不是把它们当成彼此独立的主题。
  \item 学会把 \cref{ch:polarization} 中的 little group / 不可约表示分类直接用来判断 exact $\Gamma$ 点的法向辐射允许性。
  \item 明确 BIC / quasi-BIC 在 PCSEL 中是候选模与辐射工程语言，不是绕过器件物理的捷径。
\end{itemize}

\paragraph{读前准备}
建议已掌握 \cref{ch:maxwell,ch:gamma,ch:polarization} 中关于开放系统、$\Gamma$ 点模、辐射损耗和偏振/远场的内容。

\paragraph{阅读主线}
\ChapterModelLevel{L1→L2}
PCSEL 的目标模通常位于 $\Gamma$ 点附近；设计上既需要较低辐射损耗以降低阈值增益，也需要保留可控的法向出射以形成可用远场。本章用 BIC / quasi-BIC 语言统一描述这两个要求。先定义开放周期系统中的 BIC，再把 \cref{ch:polarization} 的 little group / 不可约表示工具用于法向辐射通道，说明 exact $\Gamma$ 点某些 irrep 的 direct radiation 为何可被对称性禁止。随后讨论小对称性破缺如何打开 quasi-BIC 辐射，并说明真实 PCSEL 通常工作在 quasi-BIC 设计区，而不是无限大、完美无损的理想 BIC 点。

\section{为什么 BIC 不是旁枝话题，而是 PCSEL 主线中的辐射语言}
PCSEL 是开放系统，因此需要同时描述本征模及其向连续辐射通道的泄漏。BIC 描述的是一种\textbf{频率落在辐射连续谱内、却因干涉或对称性而不向外泄漏的极限状态}。也就是说，BIC 是在辐射区内使净辐射振幅为零，而非频率离开辐射区。

这里说的“辐射连续谱”是 slab 外自由空间中、位于光锥以内的一整族可传播平面波模，而非抽象的谱论术语。对 PCSEL 读者来说，把它先理解成“所有允许向上、向下逃逸的远场平面波通道的连续集合”就足够了。

这一区分对 PCSEL 很重要：设计目标是判断\textbf{哪些 $\Gamma$ 点候选模因对称性或干涉而辐射受限，哪些模更容易向法向出光}，而非单纯寻找高 $Q$ 模。BIC / quasi-BIC 语言把高 $Q$、偏振、法向远场和对称性归入同一辐射机制。

\begin{IntuitionBox}
对 PCSEL 来说，BIC 是描述辐射抑制及其受控打开的极限参照。实际设计通常不工作在理想 BIC 点本身，而是在其附近的 quasi-BIC 区域选择辐射损耗与输出耦合的折中。
\end{IntuitionBox}

\section{从辐射振幅到 BIC：什么叫“连续谱中的束缚态”}
在进入公式之前，先把 BIC、普通束缚态和准 BIC 的边界一次说清。普通束缚态的“不辐射”来自频率本身不在允许传播的连续谱里；BIC 的“不辐射”则更微妙：它的频率已经落在连续谱内，按频率和动量本应可以向外辐射，但所有可用辐射通道的净耦合振幅恰好为零。准 BIC 则是理想 BIC 在有限尺寸、轻微几何扰动、上下层栈不对称或工艺误差作用下打开一条很弱辐射通道后的现实版本。

把某个 $\Gamma$ 点附近候选模与自由空间辐射通道之间的耦合写成一个最小时间耦合模模型。设腔内模振幅为 $a$，上方和下方辐射通道的出射振幅记为 $s_{\uparrow,-}$ 与 $s_{\downarrow,-}$，则可写
\begin{equation}
\dot a
=
\left(\ii\omega_0-\gamma_{\mathrm{nr}}-\gamma_{\mathrm{rad}}\right)a,
\qquad
\begin{bmatrix}
s_{\uparrow,-}\\
s_{\downarrow,-}
\end{bmatrix}
=
\begin{bmatrix}
d_\uparrow\\
d_\downarrow
\end{bmatrix} a,
\label{eq:bic_tcmt}
\end{equation}
其中 $\gamma_{\mathrm{nr}}$ 表示非辐射损耗，$\gamma_{\mathrm{rad}}$ 表示辐射损耗，$d_\uparrow,d_\downarrow$ 是模到辐射通道的耦合振幅。若在某个对称性或干涉条件下，
\begin{equation}
d_\uparrow = d_\downarrow = 0,
\label{eq:bic_condition}
\end{equation}
则该模即使频率位于辐射连续谱中，也不会向外辐射；理想无损情况下，$\gamma_{\mathrm{rad}}=0$，这就是 BIC 的最小数学图像。

需要特别强调的是，\cref{eq:bic_condition} 表示\textbf{辐射振幅为零}，而不是模式根本不与外界有任何关系。只要结构轻微破坏对称性、有限尺寸边缘打开旁路通道、表面粗糙引入散射，或者工作点变化导致场型重分布，这个零就会被抬起，理想 BIC 便退化成具有有限辐射线宽的 quasi-BIC。

BIC 的定义依赖开放边界与连续辐射通道。若把问题错误地近似成封闭本征值问题，就既不会看到真正的辐射损耗，也不会正确看到 BIC 的成因。

\textbf{用干涉理解 BIC。}BIC 并不要求辐射通道不存在，而是要求通道中的净辐射振幅相消为零。对称性保护型 BIC 中，这种相消由结构对称性强制给出，而不是数值上的偶然抵消。

\begin{ExampleBox}
\textbf{BIC 的严格口径是什么？}

普通束缚态可概括为"频率不在连续谱内，所以无辐射"；BIC 则应表述为"频率位于连续谱内，但辐射耦合矩阵元为零，所以无辐射"。对 PCSEL 来说，这个区分非常关键，因为 BIC / quasi-BIC 讨论的是开放辐射问题本身，而不是简单靠频率躲开辐射区的模式。
\end{ExampleBox}

\section{对称性保护 BIC 的物理图像：从偶极子到点群表示}
\label{sec:symmetry-bic-intuitive}
第\cref{ch:polarization}已经从点群不可约表示角度给出了分类规则。本节用一个偶极子辐射例子说明，对称性为何能够使辐射矩阵元为零。

考虑局域源的多极展开，最低阶项为电偶极辐射。电偶极矩定义为
\begin{equation}
\bm{d} = \int_V \bm{P}(\bm{r}) \dd V,
\label{eq:dipole_moment}
\end{equation}
其中$\bm{P}$是极化强度分布。关键问题在于：$\bm{d}$在空间反演操作$\mathcal{I}:\bm{r}\to-\bm{r}$下如何变换？

若模式电场满足奇宇称条件$\bm{E}(-\bm{r}) = -\bm{E}(\bm{r})$，则极化分布也具有相同宇称$\bm{P}(-\bm{r}) = -\bm{P}(\bm{r})$。此时对整个单胞积分：
\begin{equation}
\bm{d} = \int_{\rm cell} \bm{P}(\bm{r}) \dd V
      = \int_{\rm cell} \bm{P}(-\bm{r}') \dd V'
      = -\int_{\rm cell} \bm{P}(\bm{r}') \dd V'
      = -\bm{d},
\end{equation}
其中用了变量代换$\bm{r}'=-\bm{r}$且积分区域不变（周期单胞关于原点对称）。于是得到$\bm{d} = -\bm{d}$，唯一可能是$\bm{d}=\bm{0}$。\textbf{电偶极矩为零意味着最低阶辐射通道被关闭}。

\begin{NoteBox}{从偶极子禁戒到点群 irrep}
上面的论证用的是最简单的空间反演对称性（宇称）。对 PCSEL 常用的正方晶格$C_{4v}$点群，情况更丰富：
\begin{itemize}
  \item $A_1,A_2$表示是一维的，场分布在某些操作下保持不变或变号
  \item $B_1,B_2$也是一维，但变换规则不同
  \item $E$表示是二维的，对应简并对
\end{itemize}
法向辐射的自由空间平面波构成$C_{4v}$的某个表示。若模式的 irrep 与辐射场的 irrep 不匹配（直积不含恒等表示），则耦合矩阵元必为零。这就是\cref{ch:polarization}中表格的物理来源。
\end{NoteBox}

将该论证推广到 $C_{4v}$ 点群：exact $\Gamma$ 点模式可按 $C_{4v}$ 不可约表示分类，法向辐射通道（沿 $z$ 轴传播的平面波）也有确定的变换行为。若模式 irrep 与辐射通道 irrep 不相容，则由群论正交性可知其重叠积分为零。因此，某些 $\Gamma$ 点模式可成为 BIC；其原因是对称性选择定则，而非数值巧合。

\section{\texorpdfstring{$\Gamma$}{Gamma} 点与 BIC 语言的对应关系}
在二维周期 slab 中，$\Gamma$ 点对应零平面波矢，恰好也是法向出射和法向辐射最直接相关的位置。对 PCSEL 来说，这种联系尤为自然：目标模式大多围绕 $\Gamma$ 点组织，因为它最容易同时连接\textbf{面内布拉格反馈}与\textbf{法向面发射}。

而在 $\Gamma$ 点，晶格对称性和场分量奇偶性经常会直接决定某一模是否能够与法向辐射通道耦合。若一个模式的近场对称性与外部法向平面波的表示不匹配，那么即使频率处于辐射连续谱中，它也可以因对称性而成为 symmetry-protected BIC。反过来，若通过轻微单胞扰动、双晶格设计或椭圆孔设计打破这类严格匹配条件，就会把原本封闭的辐射通道重新打开，得到 quasi-BIC。

因此，$\Gamma$ 点、模简并、偏振选择和法向远场不是彼此独立的主题；在 $\Gamma$ 点，它们由同一组对称性与辐射通道规则共同约束。

\section{把法向辐射通道写成 irrep：BIC 判断才真正可操作}
\cref{ch:polarization} 已给出最小分类工具：方格晶格的 $\Gamma$ 点 little group 为 $C_{4v}$，三角晶格为 $C_{6v}$。BIC 初筛还需要补充一条规则：\textbf{法向自由空间辐射通道也必须按同一 little group 分类。}

对 exact normal incidence，外部自由空间只提供两支横向平面波偏振，因此：
\begin{itemize}
  \item 在 $C_{4v}$ 下，法向辐射通道最小地张成二维表示 $E$；
  \item 在 $C_{6v}$ 下，法向辐射通道最小地张成二维表示 $E_1$。
\end{itemize}
于是，若某个 $\Gamma$ 点模的 irrep 与这些辐射通道表示不匹配，那么 direct normal-radiation matrix element 在保持该对称性的前提下就必须为零。这就是 symmetry-protected BIC 最简洁、也最适合教学复用的判断规则。

\begin{table}[htbp]
  \centering
  \footnotesize
  \caption{exact $\Gamma$ 点法向辐射允许性的一阶 irrep 判断规则。这里给出的是最小 in-plane little-group 规则；更严格的三维 slab 判断还需再结合纵向奇偶性与层栈对称。}
  \label{tab:gamma-rad-irrep}
  \setlength{\tabcolsep}{4pt}
  \begin{tabularx}{\textwidth}{@{}p{1.8cm}p{1.8cm}p{2.4cm}Y@{}}
    \toprule
    晶格 & $\Gamma$ 点 little group & 法向辐射通道 irrep & 最小结论 \\
    \midrule
    方格晶格 & $C_{4v}$ & $E$ & 属于 $E$ 的模式在对称性上允许 direct normal radiation；属于 $A_1,A_2,B_1,B_2$ 的模式可形成 symmetry-protected dark family \\
    三角晶格 & $C_{6v}$ & $E_1$ & 属于 $E_1$ 的模式在对称性上允许 direct normal radiation；属于 $A_1,A_2,B_1,B_2,E_2$ 的模式在 exact $\Gamma$ 点可成为 symmetry-protected dark family \\
    \bottomrule
  \end{tabularx}
\end{table}

这张表把“看场图判断是否辐射”的经验判断改写为可执行流程：先确定 little group，再标注 irrep，最后与辐射通道 irrep 比较。若不匹配，可优先把该模视为 symmetry-protected BIC family 候选；若匹配，再计算耦合强度。

\begin{PitfallBox}
\textbf{in-plane little-group 匹配只是必要的第一层筛选，不是三维 slab 中法向辐射强度的充分判据。} 也就是说，irrep 匹配说明 direct radiation 在对称性上“可以非零”，并不保证它一定很强；而真正的法向辐射强度还会继续受纵向奇偶性、层栈不对称、有限厚度矢量混合和有限尺寸边缘共同影响。
\end{PitfallBox}

若想把“为什么不匹配就一定不辐射”再往前推进半步，可以用最小偶极矩积分来理解。对 exact $\Gamma$ 点法向直接辐射，主导耦合常可压缩为单胞内某个零级辐射矩
\begin{equation}
\bm{d}_0 \propto \int_{V_{\mathrm{cell}}}\bm{P}_{\parallel}(\rvec)\,\dd V,
\label{eq:gamma_dipole_integral}
\end{equation}
其中 $\bm{P}_{\parallel}$ 是与法向辐射通道耦合的平面内极化或等效辐射源分量。若某个对称操作 $g$ 保持单胞与法向辐射通道不变，而模式在该操作下属于特征值 $\chi_\alpha(g)\neq 1$ 的 irrep，则
\[
\bm d_0
=
\int \bm P_{\parallel}\,\dd V
=
\int \hat T(g)\bm P_{\parallel}\,\dd V
=
\chi_\alpha(g)\bm d_0.
\]
于是只有 $\bm d_0=\bm 0$ 才可能满足上式。这给出 symmetry-protected BIC 的最小推导：该模并非只是场图较暗，而是与法向均匀辐射通道耦合的零级积分在对称性上为零。

这里还应把“辐射通道的表示”说得更正式一点。完整的法向辐射\textbf{二维偏振子空间}
在方格晶格下按 $E$ 表示变换；但当我们把耦合振幅进一步投影到某一个固定的法向线偏振
端口，或等价地研究某个固定的零级辐射矩分量时，对应的\textbf{标量重叠泛函}在保持该端口
不变的对称子群下就是平凡表示。此时若模式属于非平凡 irrep，群平均投影
\begin{equation}
\Pi_{\mathrm{rad}} =
\frac{1}{|G|}\sum_{g\in G}\chi_{\mathrm{rad}}^\ast(g)\,\hat T(g)
\label{eq:bic_group_projector}
\end{equation}
便会把该重叠压成零。上文用单个群元写出的
$\chi_\alpha(g)\neq 1\Rightarrow \bm d_0=0$，只是这个更一般投影结论的最小判据：
当存在某个保持辐射端口不变的群元，而模式在其下取得不同字符时，禁耦合就已经成立。

若某个方格晶格 slab 在理想上下对称、单胞高对称的条件下支持一个关于镜面对称为奇的 $\Gamma$ 点模，而法向自由空间平面波关于同一镜面为偶，那么该模对法向辐射的耦合振幅便可严格为零。此时它在无限大、理想结构中表现为 symmetry-protected BIC。一旦将圆孔轻微拉成椭圆，或引入双晶格扰动，这个严格零耦合被打破，模式会转化为 quasi-BIC，并出现有限但可控的法向辐射。

\section{从 BIC到 quasi-BIC：为什么真实PCSEL 更关心后者}
若只追求无穷大的被动 $Q$，理想 BIC 可以作为极限点；但 PCSEL 作为激光器还需要可控的垂直辐射，用于输出和远场整形。因此，真实器件通常选择\textbf{靠近 BIC 的 quasi-BIC 区域}，在低辐射损耗和有限输出耦合之间折中。

\cref{fig:ch08b_bic_mechanism}给出了BIC形成机制的完整可视化：(a)对称性保护BIC，当模式的不可约表示与辐射通道的表示不匹配时，耦合矩阵元严格为零；(b)Friedrich-Wintgen型BIC，两个共振模通过共享辐射连续谱发生干涉相消；(c)准BIC区域，轻微对称性破缺将理想BIC转化为具有超高但有限Q值的共振态。

\begin{figure}[htbp]
  \centering
  % !TEX root = ../main.tex
% Figure content only: inserted inside an outer figure environment in ch08b
\begin{tikzpicture}[x=1cm,y=1cm,>=Latex,font=\small]
  \tikzset{
    panel/.style={draw=gray!35,rounded corners=2pt,fill=gray!3},
    slab/.style={draw=gray!65,fill=gray!14},
    port/.style={draw=teal!65!black,fill=teal!6,rounded corners=2pt,align=center,font=\scriptsize},
    mode/.style={draw=gray!60,fill=yellow!15,circle,minimum size=0.62cm,align=center,font=\scriptsize},
    rad/.style={->,very thick,teal!70!black},
    leak/.style={->,thick,orange!90!black},
    note/.style={align=center,font=\scriptsize,text width=3.7cm}
  }

  % (a) symmetry-protected BIC
  \begin{scope}[shift={(-5.2,0)}]
    \draw[panel] (-2.3,-2.05) rectangle (2.3,2.45);
    \node[font=\bfseries] at (0,2.12) {(a) 对称性禁耦合};
    \draw[slab] (-1.8,-0.12) rectangle (1.8,0.22);
    \foreach \x in {-1.35,-0.9,...,1.35} {
      \fill[white] (\x,0.05) circle (0.08);
      \draw[gray!55] (\x,0.05) circle (0.08);
    }
    \fill[red!28] (-0.85,0.48) ellipse (0.52 and 0.20);
    \fill[blue!25] (0.85,0.48) ellipse (0.52 and 0.20);
    \node[font=\scriptsize] at (0,0.86){奇/暗模};
    \node[port] at (0,1.45){法向辐射通道};
    \draw[rad,densely dashed] (-0.45,0.66)--(-0.45,1.20);
    \draw[rad,densely dashed] (0.45,0.66)--(0.45,1.20);
    \draw[red!75!black,very thick] (-0.18,1.15)--(0.18,1.51);
    \draw[red!75!black,very thick] (-0.18,1.51)--(0.18,1.15);
    \node[note] at (0,-1.22)
      {模式表示与法向通道不匹配，辐射矩阵元被对称性压成零。};
  \end{scope}

  % (b) Friedrich-Wintgen interference
  \begin{scope}
    \draw[panel] (-2.3,-2.05) rectangle (2.3,2.45);
    \node[font=\bfseries] at (0,2.12) {(b) 干涉型 BIC};
    \node[mode,fill=orange!18] (m1) at (-0.82,0.82){模 1};
    \node[mode,fill=orange!18] (m2) at (0.82,0.82){模 2};
    \node[port,minimum width=2.6cm] (cont) at (0,-0.35){共享辐射连续谱};
    \draw[leak] (m1)--node[left,font=\scriptsize] {$+$} (cont);
    \draw[leak] (m2)--node[right,font=\scriptsize] {$-$} (cont);
    \draw[<->,thick,purple!70!black] (m1)--node[above,font=\scriptsize]{近场耦合}(m2);
    \node[draw=purple!65!black,fill=purple!6,rounded corners=2pt,font=\scriptsize,align=center]
      at (0,-1.10) {$d_1+d_2=0$};
    \node[note] at (0,-1.68)
      {两个共振模向同一连续谱泄漏，某个本征组合发生相消。};
  \end{scope}

  % (c) useful quasi-BIC window
  \begin{scope}[shift={(5.2,0)}]
    \draw[panel] (-2.3,-2.05) rectangle (2.3,2.45);
    \node[font=\bfseries] at (0,2.12) {(c) quasi-BIC 工作区};
    \draw[->,thick] (-1.55,-1.18)--(1.75,-1.18) node[right]{$\delta$};
    \draw[->,thick] (-1.55,-1.18)--(-1.55,1.40) node[above]{$Q_{\mathrm{rad}}$};
    \draw[red!75!black,very thick,domain=-1.18:1.48,samples=120]
      plot(\x,{1.8*exp(-1.2*(\x+1.18))-0.55});
    \draw[dashed,gray!65] (-1.18,-1.18)--(-1.18,1.22);
    \node[font=\scriptsize,gray!70,rotate=90] at (-1.37,0.22){理想 BIC};
    \draw[green!60!black,thick,densely dashed] (-0.18,-0.90) rectangle (1.05,-0.12);
    \node[font=\scriptsize,green!45!black,align=center] at (0.43,-0.50)
      {有限出光\\可控损耗};
    \draw[rad] (0.70,-0.08)--(0.70,0.62);
    \node[font=\scriptsize] at (0.32,-1.58){$Q_{\mathrm{rad}}\propto\delta^{-2}$};
  \end{scope}
\end{tikzpicture}

  \caption{BIC 形成机制示意图。(a) 对称保护 BIC；(b) 干涉型 BIC；(c) quasi-BIC 区域及其 $Q\propto\delta^{-2}$ 标度。}
  \label{fig:ch08b_bic_mechanism}
\end{figure}

\cref{fig:ch08b_bic_mechanism}概括了 BIC 的两类形成机制：对称性禁耦合和共用连续谱中的干涉相消。对 PCSEL 设计而言，问题是选择 quasi-BIC 区域，使辐射 $Q$、增益、阈值和远场需求相互匹配，而非简单追求 BIC。

在最简单的对称破缺参数 $\delta$ 下，辐射振幅常可近似线性展开为
\begin{equation}
d \propto \delta,
\qquad
\gamma_{\mathrm{rad}} \propto |d|^2 \propto \delta^2,
\qquad
Q_{\mathrm{rad}} \propto \delta^{-2}.
\label{eq:q_scaling_quasibic}
\end{equation}

\begin{NoteBox}{$Q\propto\delta^{-2}$的量纲分析与物理论证}
设对称性破缺参数为$\delta$，辐射振幅$d$是最小的线性响应量，则 Taylor 展开给出
$d = d_0 + d_1\delta + O(\delta^2)$。由于$\delta=0$时对应理想 BIC ($d_0=0$)，故主导项为$d\approx d_1\delta$。
辐射功率正比于$|d|^2$，因而$\gamma_{\rm rad}\propto\delta^2$，再由$Q=\omega_0/2\gamma_{\rm rad}$即得$Q\propto\delta^{-2}$。

从量纲角度看，$\delta$是无量纲几何扰动（如位移除以晶格常数），$Q$也是无量纲量，因此幂次关系只能是纯数。
二次方的根源在于：辐射是偶极矩（或更高阶多极矩）的平方，而小扰动下多极矩的变化是线性的。
\end{NoteBox}

\cref{eq:q_scaling_quasibic} 往往还可进一步写成更适合工程估算的形式
\begin{equation}
Q_{\mathrm{rad}} \approx \frac{Q_0}{\delta^2},
\qquad
\delta=\frac{d}{a},
\label{eq:q_delta_explicit}
\end{equation}
其中 $d$ 是破缺位移或等效几何扰动幅度，$a$ 是晶格常数，$Q_0$ 吸收了结构 family、工作带支和归一化定义相关的常数。\cref{eq:q_delta_explicit} 的直接后果是：若其余条件近似不变，把 $\delta$ 从 $0.02$ 降到 $0.01$，辐射 $Q$ 通常会提升约四倍。该二次标度在对称破缺 quasi-BIC 文献中有明确支撑\cite{koshelev2018quasibic}。它说明小几何扰动通常以连续方式打开辐射通道。设计问题因此转化为：在给定输出功率、阈值增益预算、偏振纯度和远场要求下选择 $\delta$ 的工作区间，而不是机械追求无限接近 BIC。

这也构成 quasi-BIC 工作区 与“最优单点设计”的差别。若把 $\delta$ 压得过小，也许被动 $Q$ 极高，但法向提取太弱、有限尺寸边缘和工艺扰动又会变得更敏感；若把 $\delta$ 放得过大，则目标模虽然更易出光，却可能因辐射损耗和竞争模重新排序而抬高阈值。PCSEL 设计更关心器件约束下可容差的 quasi-BIC 工作区，而不是数学上的极限点。

\section{前沿补充：quasi-BFIC 与宽角高 \texorpdfstring{$Q$}{Q} 设计}
除“单点附近的 quasi-BIC”外，近年还出现了 quasi-bound flat bands in the continuum（quasi-BFIC）路线：其目标是在更宽的角度/动量区间维持较高 $Q$ 与可用辐射耦合，而非只在单一动量点获得高 $Q$。对 PCSEL 设计的意义是两点：第一，工艺与角度容差可提高；第二，有限尺寸阵列中的角谱展宽不再必然迅速拉低整体辐射品质。该结论在 2025 年直接工作中已有证据支撑\cite{qin2025quasibfic}。
但 quasi-BFIC 是补充：前者更强调宽角鲁棒，后者更强调中心点高选择性，而非“替代”传统 $\Gamma$ 点 quasi-BIC。工程上应根据目标（窄主瓣极限、容差、功率提取）选择路线；BIC/quasi-BIC 综述与标度工作更适合作为背景理论锚点\cite{hsu2016,koshelev2018quasibic}。
进一步说，quasi-BFIC 的“宽角高 $Q$”结论本身仍是开放周期或大尺寸近似下的结论；一旦转入有限阵列，边缘散射与角谱离散化会显著改写可用窗口。因此，任何“宽角鲁棒”结论都必须补交有限尺寸证据（阵列尺寸、边界模型、收敛曲线与误差样本），否则只能算机理提示而非器件级判据\cite{wang2025finitesizepcsel}。

\CrossPlatformTriad
{BIC/quasi-BIC 理论锚点与 quasi-BFIC 直接证据联合使用\cite{hsu2016,koshelev2018quasibic,qin2025quasibfic}}
{需要在同一器件目标下同时给出角谱范围、有限尺寸边界、远场与模排序收敛}
{不可把“宽角高 $Q$”直接等价为“器件阈值更低或更鲁棒”；需经过有限阵列与电热回写验证}

\section{进阶：far-field polarization vortex 与拓扑电荷语言}
前文已经建立 BIC 的最小对称性与干涉描述。现代 BIC 理论还常使用动量空间远场偏振矢量场的奇点来描述 BIC，即 far-field polarization vortex。

为说明这种偏振绕转，把某一支 slab 共振模在上方辐射半空间的远场振幅写成二维复向量
\begin{equation}
\bm{c}(\kvec_\parallel)
=
\begin{bmatrix}
c_x(\kvec_\parallel)\\
c_y(\kvec_\parallel)
\end{bmatrix},
\label{eq:bic_farfield_vector}
\end{equation}
其中 $\kvec_\parallel=(k_x,k_y)$ 是平面内波矢，$c_x,c_y$ 表示该模向两支法向线偏振通道的辐射振幅。BIC 点的最小判据就是
\begin{equation}
\bm{c}(\kvec_{\parallel,\mathrm{BIC}})=\bm{0}.
\label{eq:bic_farfield_zero}
\end{equation}
也就是说，BIC 是远场矢量场中的零点。

一旦离开这个零点，远场会重新出现，并携带某个主偏振方向。对近线偏振远场，主轴角可先写成
\begin{equation}
\phi(\kvec_\parallel)
=
\frac{1}{2}\arg\!\left[c_x^2(\kvec_\parallel)+c_y^2(\kvec_\parallel)\right],
\label{eq:bic_polarization_angle}
\end{equation}
它把 \cref{eq:bic_farfield_vector} 中的复振幅 $c_x,c_y$ 直接转成偏振椭圆长轴偏角。于是，若沿着包围 BIC 的闭合回路 $C$ 绕一圈，主偏振方向角 $\phi(\kvec_\parallel)$ 发生整数倍 $2\pi$ 的绕转，则可定义拓扑电荷
\begin{equation}
q
=
\frac{1}{2\pi}\oint_C \nabla_{\kvec_\parallel}\phi\cdot \dd \kvec_\parallel.
\label{eq:bic_topological_charge}
\end{equation}
这个整数 $q$ 就是 far-field polarization vortex 的 winding number。它把“这个 BIC 为什么不是数值偶然出现的零点”转化成了一个拓扑语言：只要相关带支与远场矢量场连续变化，这个零点不能凭空消失，而必须通过\textbf{带着总电荷守恒地移动、成对产生或成对湮灭}的方式演化\cite{zhen2014topological,hsu2016,yang2014bic}。

\begin{figure}[htbp]
  \centering
  \begin{tikzpicture}[scale=1.0,>=Latex]
    \draw[->] (-2.2,0)--(2.3,0) node[right]{$k_x$};
    \draw[->] (0,-2.2)--(0,2.3) node[above]{$k_y$};
    \draw[dashed,gray] (0,0) circle (1.45);
    \node[gray] at (1.85,1.55){$C$};
    \filldraw[fill=black,draw=black] (0,0) circle (1.5pt);
    \node[above right] at (0,0){连续谱束缚态，$q=+1$};
    \foreach \ang in {0,30,...,330}{
      \pgfmathsetmacro{\x}{1.45*cos(\ang)}
      \pgfmathsetmacro{\y}{1.45*sin(\ang)}
      \pgfmathsetmacro{\tx}{0.45*cos(\ang+90)}
      \pgfmathsetmacro{\ty}{0.45*sin(\ang+90)}
      \draw[blue!70!black,->,line width=0.8pt] (\x,\y)--(\x+\tx,\y+\ty);
    }
  \end{tikzpicture}
  \caption{$k$ 空间中远场偏振涡旋的最小示意。闭合回路 $C$ 围绕连续谱束缚态奇点一周时，远场主偏振方向同步绕转一周，对应拓扑电荷 $q=+1$。}
  \label{fig:bic_polarization_vortex}
\end{figure}

拓扑电荷保护的对象，是\emph{动量空间远场矢量场中的奇点结构}，不是器件性能的一切性质。也就是说，电荷守恒可以帮助我们追踪 BIC 在 $\kvec_\parallel$ 空间和参数空间中的迁移、生成与湮灭，但它并不自动保证最低阈值、最佳热稳定性或最佳注入均匀性。对 PCSEL 而言，拓扑语言是用来追踪候选模辐射家族的，不是用来跳过器件模型的。

这套语言对 $\Gamma$ 点附近模式尤其有用。对称性保护型 $\Gamma$ 点 BIC 可以理解成\textbf{被晶格对称性钉扎在区中心的偏振涡旋中心}。当单胞扰动、双晶格参数或外延不对称连续变化时，读者不必只盯着“$Q$ 有没有升高”，还可以追问：
\begin{itemize}
  \item 这个偏振涡旋中心是否仍停留在 $\Gamma$ 点，还是已经移到 off-$\Gamma$；
  \item 总拓扑电荷是否守恒，是否出现了成对移动的 BIC 奇点；
  \item 目标模式是仍属于原来的 quasi-BIC family，还是已经越过了亮模/暗模家族的分界。
\end{itemize}

现代 BIC 拓扑语言在这里的作用，是提供一种比单独比较谱线宽度更稳定的模追踪方法。特别在 $\Gamma$ 点附近模式接近、偏振与远场同时变化时，far-field polarization vortex 与拓扑电荷可辅助判断模族是否连续演化。

设某个方格晶格 PCSEL 单胞在高对称参数处支持一个位于 $\Gamma$ 点的 symmetry-protected BIC。随着双晶格扰动参数缓慢增大，直接法向辐射被逐步打开，目标模退化为 quasi-BIC。此时若只看单点 $Q$ 值，读者可能只会得到“$Q$ 在下降”这一条信息；而若同时追踪 $\bm{c}(\kvec_\parallel)$ 的零点与绕转方向，就能看出该模究竟是在沿同一拓扑家族连续偏离 BIC，还是已经与邻近亮模发生重组。前者更适合解释为“可控打开辐射”，后者则意味着模式排序、偏振和远场可能一起改写。

\section{对称性破缺在 irrep 语言下如何打开 quasi-BIC}
使用 little group 工具后，可以更正式地表述“对称性破缺打开 quasi-BIC”。设高对称结构在 $\Gamma$ 点支持某个 dark family，其 irrep 与法向辐射通道 irrep 不匹配，因此 exact $\Gamma$ 点 direct radiation 为零。若小扰动把点群降低到某个子群，则原 irrep 会在子群下重新分解。常见结果有两类：
\begin{enumerate}
  \item 原本的二维 irrep 分裂成两个 singlet，于是偏振被钉扎、频率与损耗同时分裂；
  \item 原本与辐射通道不匹配的 dark family，在子群中与辐射通道落入同一表示，于是零耦合被抬起，BIC 退化为 quasi-BIC。
\end{enumerate}
这一步把 \cref{ch:polarization} 的偏振选择/模式分裂与本章的 BIC 打开/关闭联系起来。二者都是降群问题的结果：前者关注本征频率与偏振，后者关注辐射矩阵元与远场。

\section{BIC 如何连接偏振选择与远场工程}
\cref{ch:polarization} 已经说明，偏振与远场由近场矢量结构和辐射通道共同决定。用 BIC 语言表述，\textbf{某一偏振能否辐射、沿何方向辐射以及远场中心亮暗，取决于该模与相应 BIC 条件的距离。}

对称保护型 BIC 常与特定偏振表示绑定。更严格地说，在方格晶格里这常表现为某一支 $E$ 双重态在降群后更强地获得与法向辐射通道的匹配，而另一支仍更接近 dark family；在三角晶格里则常表现为 $E_1$ 和 $E_2$ family 在小扰动下偏离 BIC 条件的速度不同。于是，当设计者通过单胞形变打破旋转对称性时，实际上做的并不只是“让偏振分裂”，更是在改变不同偏振支路偏离 BIC 条件的速度。某一偏振支路若仍更接近辐射抑制条件，其辐射损耗会更低、阈值增益更小；另一偏振支路若更远离该条件，则更容易辐射、更容易在远场中形成主瓣，但被动损耗也更高。

从远场角度看，若某模严格处于 symmetry-protected BIC 点，它在法向中心常表现为零辐射；小扰动打破对称性后，法向中心会逐步亮起，远场主瓣和偏振纯度随之变化。因此，许多远场整形方法可表述为控制目标模式偏离 BIC 条件的方式。

把 BIC 理解成“理想 PCSEL 就应该完全不辐射”，是概念错误。激光器若完全不向目标通道辐射，就无法作为可用面发射器件工作。BIC 在 PCSEL 中的作用是帮助设计者理解辐射如何被抑制与重开，而不是让器件退化成只会储能的封闭态。

\section{为什么 BIC 仍然不能替代阈值与器件模型}
到这里必须再次把层级边界说清楚。BIC / quasi-BIC 语言擅长回答的问题是：\textbf{在开放周期光学层面，辐射损耗为什么高或低、怎样通过对称性与干涉调控它、以及这如何影响偏振与远场。}但它不能直接回答阈值电流、热稳定性和真实工作线宽。

原因很直接。首先，理想 BIC 讨论的是无限大、精确对称、无工艺误差的被动结构；真实器件里存在有限阵列边缘、粗糙、刻蚀误差、金属和注入非均匀。其次，即使某个模式具有非常小的辐射损耗，也未必与量子阱重叠最好，未必避开重掺杂损耗层，也未必在热回写后继续保持优势。最后，线宽和调制稳定性还涉及载流子、增益压缩和相位噪声，这已经超出了 BIC 自身的描述范围。

\begin{CommonPitfallBox}{无限周期到有限器件的外推}
\textbf{错误理解：}“无限周期仿真中该模接近 BIC，高 $Q$，所以有限器件必然高 $Q$ 且低阈值。”
\textbf{正确做法：}把无限周期结果视为候选筛选，随后必须报告有限阵列尺寸、边界截断、粗糙模型，以及 $Q$ 与远场随尺寸收敛曲线。
\textbf{最小报告项：}至少给出阵列尺寸、边界模型、收敛曲线、误差样本四要素。
\end{CommonPitfallBox}

\FiniteSizeReportTemplate

因此，本章在全书中的位置是：把前文关于 $\Gamma$ 点与偏振/远场的描述改写为开放系统辐射语言，并为后续 CWT、全波和器件章节提供候选模筛选工具，而不在本章给出器件级结论。

\section*{回看本章}
BIC / quasi-BIC 为 PCSEL 提供了连接 $\Gamma$ 点、辐射损耗、偏振和远场工程的开放系统语言。BIC 的关键是在连续谱中通过对称性或干涉使辐射振幅为零，而非频率离开辐射连续谱；在 $\Gamma$ 点，这可通过比较模式 irrep 与法向辐射通道 irrep 来判断。若二者不匹配，则 direct radiation 在对称性上被禁止。quasi-BIC 则是在真实设计中通过小而可控的对称破缺重新打开辐射通道。PCSEL 通常关注理想 BIC 附近的工作区，以同时兼顾低辐射损耗、可用面发射、偏振选择和远场工程；BIC 是阈值与器件模型的前端语言，但不能替代后续器件分析。

\section*{练习题}
\begin{enumerate}
  \item \basicq 用自己的话解释为什么 BIC 的定义必须建立在开放边界与辐射连续谱上，而不能在封闭腔语言中直接照搬。

  \item \basicq 结合 \cref{eq:bic_condition} 说明“耦合到连续谱”与“向连续谱辐射”为何不是同一句话。

  \item \midq 对方格晶格，说明为什么 exact $\Gamma$ 点法向辐射通道属于 $E$，并据此解释为什么一维 irrep family 更容易成为 symmetry-protected BIC 候选。

  \item \midq 利用 \cref{eq:q_scaling_quasibic} 讨论为什么 quasi-BIC 设计常出现“越靠近理想 BIC，越需要警惕工艺敏感性”的权衡。

  \item \hardq 选择一种你熟悉的单胞扰动方式，说明它更可能首先改写的是偏振分裂、法向远场，还是辐射损耗，并解释原因。
\end{enumerate}

\section*{延伸阅读}
\begin{itemize}
  \item 下一章开始进入 CWT 语言。本章的 irrep 与辐射选择规则将在那里对应到耦合矩阵中的损耗项和模式分裂。

  \item BIC 的一般背景可参考 \cite{fan2002guided,hsu2013,hsu2016,kodigala2017}；far-field polarization vortex 与拓扑电荷可参考 \cite{zhen2014topological,yang2014bic}。

  \item 阅读后续器件章节时，应检查 quasi-BIC 设计在有限尺寸、注入非均匀和热回写后是否仍保持目标模优势。
\end{itemize}
